\documentclass[dvipsnames]{article}

\usepackage[utf8]{inputenc}
\usepackage{fontspec}
\usepackage{polyglossia}
\usepackage{unicode-math}
\usepackage{csquotes}
\setdefaultlanguage{russian}
\setotherlanguage{english}

\setmainfont{CMU Serif}
\setsansfont{CMU Sans Serif}
\setmonofont{CMU Typewriter Text}
% \setmathfont{Latin Modern Math}

\usepackage[letterpaper,top=2cm,bottom=2cm,left=3cm,right=3cm,marginparwidth=1.75cm]{geometry}

\usepackage{amsmath}
\usepackage{graphicx}
\usepackage{amsthm}
\usepackage{array}
\usepackage{fancyhdr}
\usepackage{bussproofs}
\usepackage[colorlinks=true, allcolors=blue]{hyperref}
\usepackage[backend=biber]{biblatex}
\usepackage{xcolor}

\usepackage{minted2}
\usemintedstyle{xcode}
\setminted{escapeinside=||}
\newminted[stella]{cpp}{escapeinside=||,frame=single}
\newmintinline[haskell]{haskell}{fontsize=\scriptsize}
\newmintinline[java]{java}{fontsize=\scriptsize}
\newmintinline[python]{python}{fontsize=\scriptsize}
\newmintinline[code]{c}{fontsize=\small}
\let\codeOrig\code
\renewcommand{\code}[2][]{%
  \ifmmode
    \if\relax\detokenize{#1}\relax
      \mathchoice{\mbox{\normalfont\small\ttfamily #2}}{\mbox{\normalfont\small\ttfamily #2}}{\mbox{\normalfont\small\ttfamily #2}}{\mbox{\normalfont\small\ttfamily #2}}%
    \else
      \mathchoice{\mbox{\normalfont\small\ttfamily #2}}{\mbox{\normalfont\small\ttfamily #2}}{\mbox{\normalfont\small\ttfamily #2}}{\mbox{\normalfont\small\ttfamily #2}}%
    \fi
  \else
    \codeOrig[#1]{#2}%
  \fi}

\usepackage{sourcecodepro}
\usepackage{mdframed}

\DeclareMathOperator{\xor}{\mathsf{xor}}
\DeclareMathOperator{\tru}{\mathsf{tru}}
\DeclareMathOperator{\fls}{\mathsf{fls}}

\addbibresource{../references.bib}

\begin{document}

\pagestyle{fancy}
\fancyhf{}
\chead{Задание выполнил \textcolor{blue}{Бенжамин Пирс (замените на своё имя)}}
\cfoot{Языки программирования. Семантика и системы типов (весна 2026, Университет ИТМО)}

\section*{Тема 4. Теоретическое задание \textcolor{blue}{(выполнил Бенжамин Пирс)}}%

Представьте следующий вариант обобщённого цикла \code{for}:
  \begin{enumerate}
    \item[\textbf{Синтаксис}] \quad \\
      \begin{stella}
t ::= ...
      for/result[E, T] (x=t; t; t) do t
      \end{stella}
    \item[\textbf{Семантика}] \quad \\ Цикл \code{for/result[E, T] (x=t|$_1$|; t|$_2$|; t|$_3$|) do t|$_4$|}
    инициализирует локальную переменную \code{x} типа \code{T} значением выражения \code{t|$_1$|},
    после чего следует неограниченное число итераций; на каждой итерации проверяется
    условие \code{t|$_2$|} (может содержать \code{x}); если условие верно, то
    выполняется тело цикла \code{t|$_4$|} (может содержать \code{x}), которое возвращает
    значение варианта типа \code{E + Unit}:
    либо \code{inl v|$_\mathsf{err}$|}, означающее ошибку, либо \code{inr unit}, означающее успешное продолжение.
    Если \code{t|$_4$|} на некоторой итерации возвращает \code{inl v|$_\mathsf{err}$|}, то выполнение цикла немедленно
    останавливается, и результатом всего цикла становится \code{inl v|$_\mathsf{err}$|}.
    Если же \code{t|$_4$|} на всех итерациях возвращает \code{inr unit}, то после каждого успешного выполнения тела
    выполняется обновление \code{x} через выражение \code{t|$_3$|} (может содержать \code{x}) и начинается следующая итерация.
    Когда в некоторый момент условие \code{t|$_2$|} становится ложным, цикл завершается и результатом всего цикла
    становится \code{inr x}, где \code{x} — последнее значение аккумулятора.
    Семантика этой версии цикла \code{for/result} примерно соответствует семантике следующего блока кода на языке Python:
    \begin{minted}[escapeinside=||,texcomments]{python}
x = t|$_1$|               # инициализируем переменную x
while t|$_2$|:            # пока верно условие t$_2$
  res = t|$_4$|           # выполняем тело цикла t$_4$
  match res:
    case inl(v_err):      # если произошла ошибка, немедленно выходим
      return inl(v_err)
    case inr(unit):   # иначе продолжаем цикл
      x = t|$_3$|          # обновляем x (возможно, используя текущее значение x)
return inr(x)        |$\phantom{_1}$|# возвращаем успешный результат с последним значением x
    \end{minted}
    \item[\textbf{Типизация}] \quad \\
    \begin{prooftree}
      \AxiomC{$\Gamma \vdash \code{t}_1\code{ : T}$}
      \AxiomC{$\Gamma, \code{x : T} \vdash \code{t}_2\code{ : Bool}$}
      \AxiomC{$\Gamma, \code{x : T} \vdash \code{t}_3\code{ : T}$}
      \AxiomC{$\Gamma, \code{x : T} \vdash \code{t}_4\code{ : E + Unit}$}
      \RightLabel{T-For}
      \QuaternaryInfC{$\Gamma \vdash$ \code{for/result[E, T] (x=t|$_1$|; t|$_2$|; t|$_3$|) do t|$_4$| : E + T}}
    \end{prooftree}
  \end{enumerate}

  \clearpage
\begin{enumerate}

  \item Выпишите функцию раскрытия сокращений, соответствующую такому варианту цикла \code{for/result},
  для простого типизированного $\lambda$-исчисления~\cite[Рис.~9.1]{TaPL-ru} с логическими выражениями~\cite[Рис.~8.1]{TaPL-ru},
  комбинатором неподвижной точки~\cite[Рис.~11.12]{TaPL-ru} (исключая \texttt{letrec}!)
  и вариантами~\cite[Рис.~11.11]{TaPL-ru} (типы-суммы в базовое исчисление не добавляются!).

  \color{blue}
  \textbf{Решение.}
  \emph{место для решения}
  \color{black}

  \clearpage

  \item Покажите, что функция раскрытия сокращений сохраняет типизацию, если возможно.
  Иначе — продемонстрируйте на контрпримере, почему сохранение типизации невозможно.

  \color{blue}
  \textbf{Решение.}
  \emph{место для решения}
  \color{black}

  \clearpage

  \item Можно ли убрать явную аннотацию типов \code{for/result[E, T]}? В любых выражениях или только в некоторых?
  Аргументируйте свой ответ.

  \color{blue}
  \textbf{Решение.}
  \emph{место для решения}
  \color{black}

\end{enumerate}

\section*{Список литературы}
\printbibliography[heading=none]

\end{document}

